\section{Four developed case studies}
\label{sec:casestudies}

\subsection{Exact bisection for the inverse of \texorpdfstring{$y^3+y$}{y3+y}}

Let
\[
        f(y)=y^3+y.
\]
For rational input \(x\), we construct a regular rational name for the unique
solution of \(f(y)=x\).  The construction is useful because it exercises exact
rational arithmetic, a bounded loop, a monotonicity certificate, and a final
error contract without any transcendental machinery.

\begin{lemma}[Strict monotonicity]
For all rationals \(a<b\), \(f(a)<f(b)\).
\end{lemma}

\begin{proof}
Factor the difference:
\[
 f(b)-f(a)
 =(b-a)(a^2+ab+b^2+1).
\]
The first factor is positive.  The second is positive because
\[
 a^2+ab+b^2=\left(a+\frac b2\right)^2+\frac{3b^2}{4}\ge0.
\]
All identities are rational polynomial identities and are therefore suitable
for iPRA after denominator clearing.
\end{proof}

Set \(B=|x|+1\).  Then
\[
        f(-B)\le -B\le x\le B\le f(B),
\]
so \([-B,B]\) brackets the unique solution.  Define a bisection trace.  At
state \((\ell_i,h_i)\), let \(m_i=(\ell_i+h_i)/2\).  If
\(f(m_i)\le x\), set \((\ell_{i+1},h_{i+1})=(m_i,h_i)\); otherwise set it to
\((\ell_i,m_i)\).

\begin{lemma}[Bisection invariant]
For every \(i\),
\[
 f(\ell_i)\le x\le f(h_i),
 \qquad
 h_i-\ell_i=2B\cdot2^{-i}.
\]
\end{lemma}

\begin{proof}
The base case is the bracket above.  At a step, rational comparison decides
\(f(m_i)\le x\).  In the first branch, the lower inequality is updated to
\(f(m_i)\le x\), and the upper inequality is unchanged; the second branch is
symmetric.  In either branch the interval width is halved.  The conjunction is
a fixed Boolean combination of polynomial inequalities and rational equalities,
so it can be the trace invariant.
\end{proof}

Let \(q_p\) be the midpoint after \(N(p)\) steps, where \(N(p)\) is the
least integer satisfying
\[
       2B\cdot2^{-N(p)}\le 2^{-p-2}.
\]
Bounded doubling from the numerator and denominator of \(B\) computes this
integer primitively recursively.  Taking the least such integer makes \(N\)
nondecreasing, so the intervals used for successive precisions are nested.

\begin{theorem}[Cubic inverse name]
The sequence \((q_p)\) is a regular Cauchy name.  It determines the unique
real solution of \(y^3+y=x\), and every finite assertion needed for this claim
is provable from the bisection trace in iPRA.
\end{theorem}

\begin{proof}
The solution lies in every bracketing interval.  Hence the midpoint is within
half the interval width, at most \(2^{-p-3}\), of the solution.  For
\(r>p\), both \(q_p\) and \(q_r\) are within their respective half-widths of
the same nested solution interval, so
\[
       |q_r-q_p|\le 2^{-p-2}.
\]
This is stronger than the usual regularity requirement after a harmless index
shift.  Strict monotonicity gives uniqueness.  Internally, one need not
quantify over a completed real: the theorem states nested rational brackets,
shrinking widths, and the Cauchy inequality between midpoint terms.  These are
bounded rational calculations compiled by trace reflection.
\end{proof}

The included Python certificate for \(x=2\) and precision parameter \(12\)
records every rational interval and is independently checked by exact
cross-multiplication.

\subsection{Height bounds and complete lists}

Let \(P(c)\) be a decidable predicate on coded points, let \(h(c)\) be a
primitive-recursive height, and suppose there is a primitive-recursive finite
universe \(U(H)\) containing every well-formed code of height at most \(H\).

\begin{theorem}[Height bound implies complete list]
If
\[
        \forall c\,[P(c)\to h(c)\le H],
\]
then bounded filtering of \(U(H)\) by \(P\) produces a complete finite list of
all points satisfying \(P\).
\end{theorem}

\begin{proof}
Every retained code satisfies \(P\) by the filter definition.  Conversely, if
\(P(c)\), the height assumption puts \(c\) in the represented finite universe
\(U(H)\); the filter therefore retains it.  Duplicates may be removed by
bounded equality search.  All loops are bounded by the size of \(U(H)\).
\end{proof}

\begin{theorem}[Complete list implies height bound]
Suppose a finite list \(L\) is certified sound and complete:
\[
  c\in L\to P(c),
  \qquad
  P(c)\to c\in L.
\]
Then
\[
        H=\max\{h(c):c\in L\},
\]
with \(H=0\) for the empty list, is a height bound for every \(P\)-point.
\end{theorem}

\begin{proof}
If \(P(c)\), completeness gives \(c\in L\).  The defining finite-maximum
invariant says that the maximum is at least the height of every list entry.
Therefore \(h(c)\le H\).  The maximum is a primitive-recursive fold.
\end{proof}

The equivalence is often misunderstood.  It does \emph{not} say that a proof
of abstract finiteness yields a cutoff.  It says that either of two explicit
quantitative witnesses---a height bound or a certified complete list---can be
converted to the other by finite computation.

\subsection{Choice-independent simplicial maps from carriers}

Let \(K,L\) be finite simplicial complexes.  A strong carrier assigns to each
vertex \(v\) of \(K\) a nonempty finite set \(C(v)\) of vertices of \(L\), with
the property that for every simplex \(\sigma\in K\),
\[
      \bigcup_{v\in\sigma} C(v)
\]
is contained in a simplex of \(L\).  A selection chooses one
\(f(v)\in C(v)\) for each vertex.

\begin{theorem}[Carrier independence]
Every selection is a simplicial map \(f:K\to L\).  Any two selections
\(f,g\) are contiguous, and therefore induce the same homomorphism on homology
over \(\F_2\).
\end{theorem}

\begin{proof}
For a simplex \(\sigma\), the selected vertices \(f(v)\) lie in the union of
its carrier clouds, which lies in a target simplex.  Thus \(f\) is simplicial.
For two selections, the union
\(f(\sigma)\cup g(\sigma)\) is still contained in the same union of carrier
clouds, so it is a simplex; this is contiguity.

Define the prism operator on an ordered \(k\)-simplex by
\[
 P[v_0,\ldots,v_k]
 =\sum_{i=0}^k
 [f(v_0),\ldots,f(v_i),g(v_i),\ldots,g(v_k)]
 \quad\text{over }\F_2,
\]
omitting degenerate simplices according to the fixed normalization convention.
A finite cancellation calculation gives
\[
       \partial P+P\partial=f_\#+g_\#.
\]
Hence, for a cycle \(z\),
\(f_\#z+g_\#z=\partial Pz\), so the two images represent the same homology
class.  Every chain group and sum is finite; the prism identity is a bounded
matrix equality and can be certified directly.
\end{proof}

This theorem is the correct constructive replacement for choosing a canonical
simplicial approximation.  A finite carrier cloud is enough, because the
observable homology map is choice independent.

\subsection{Strong finite affine gluing}

Let \(R\) be a coded commutative ring and \(M\) a coded \(R\)-module.  Fix
\(f_0,\ldots,f_{r-1}\in R\) and an exponent \(N\), together with a cover
certificate
\begin{equation}
        \sum_{i<r} a_i f_i^N=1.
        \label{eq:coveridentity}
\end{equation}
A local section on \(D(f_i)\) is represented by \(m_i/f_i^N\).  Assume the
strong synchronized compatibility equations
\begin{equation}
        f_j^N m_i=f_i^N m_j
        \qquad(i,j<r).
        \label{eq:strongcompat}
\end{equation}
Define
\[
             m=\sum_{i<r}a_i m_i.
\]

\begin{theorem}[Finite affine gluing]
Under \eqref{eq:coveridentity} and \eqref{eq:strongcompat}, the global element
\(m\in M\) restricts to \(m_k/f_k^N\) on every \(D(f_k)\).  It is unique among
global elements whose differences are annihilated by a supplied common cover
power.
\end{theorem}

\begin{proof}
In the localization at \(f_k\), use the cover identity:
\[
 \frac{m_k}{f_k^N}
 =\sum_{i<r}a_i f_i^N\frac{m_k}{f_k^N}
 =\sum_{i<r}a_i m_i
 =m.
\]
The middle equality is exactly the localized form of
\eqref{eq:strongcompat}.  For uniqueness, if \(m'\) has the same restrictions
and a common exponent \(L\) satisfies
\(f_i^L(m-m')=0\) for every \(i\), use a certified identity
\(\sum b_i f_i^L=1\) to obtain
\[
 m-m'=\sum b_i f_i^L(m-m')=0.
\]
All sums, powers, and module equations are finite.
\end{proof}

Ordinary localization compatibility may require increasing denominators and
finding a common annihilating exponent.  The framework treats that exponent
and the associated cover identity as a finite producer certificate.  Once it
is supplied, gluing is a bounded algebraic trace.  This yields a concrete entry
point for affine schemes and coherent sheaves in iPRA without pretending that
arbitrary localization witnesses have a uniform bound for free.
